\section{Coordinate Changes}

\begin{exercise}
    A set $V \subset \A^n(k)$ is called a \textit{linear subvariety} of $\A^n(k)$ if $V = V(F_1, ..., F_r)$ for some polynomials $F_i$ of degree 1. (a) Show that if $T$ is an affine change of coordinates on $\A^n$, then $V^T$ is also a linear subvariety of $\A^n(k)$. (b) If $V \neq \varnothing$, show that there is an affine change of coordinates $T$ of $\A^n$ such that $V^T = V(X_{m+1}, ..., X_n)$. (\textit{Hint:} use induction on $r$.) So $V$ is a variety. (c) Show that the $m$ that appears in part (b) is independent of the choice of $T$. It is called the dimension of $V$. Then $V$ is then isomorphic (as a variety) to $\A^m(k)$. (\textit{Hint:} Suppose there were an affine change of coordinates $T$ such that $V(X_{m+1}, ..., X_n)^T = V(X_{s+1}, ..., X_n)$, $m < s$; show that $T_{m+1}, ..., T_n$ would be dependent.) \\
\end{exercise}

\begin{solution}
    \begin{enumerate}[label=(\alph*)]
        \item If we denote $I = (F_1, ..., F_r)$, then $I^T$ is generated by the elements $F^T$ with $F \in I$. But notice that if $F = \sum_i a_iF_i$, then
        $$F^T = \tilde{T}(\sum_i a_iF_i) = \sum_i \tilde{T}(a_i)\tilde{T}(F_i) = \sum_i a_i^T F_i^T$$
        which implies that $I^T = (F_1^T, ..., F_r^T)$. Since the $F_i$'s and the $T_i$'s are all linear, then the $F_i^T$ are linear, and hence, $V^T = V(F_1^T, ..., F_r^T)$ is a linear subvariety.
        \item Let's prove it by induction on $r\geq 1$. When $r = 1$, we have that $V = V(f)$ where $\deg(f) = 1$. If we write $f = a_0 + \sum_i a_iX_i$, then we must have at least one $i$ such that $a_i \neq 0$, call it $i = i_0$. Define the polynomial map $T = (T_1, ..., T_n)$ where $T_i = 0$ for all $i \neq i_0$, and $T_{i_0} = a_{i_0}^{-1}(X_n - a_0)$, then
        $$V^T = V(f^T) = V\left(a_0 + \sum_i a_iT_i\right) = V\left(a_0 + a_{i_0}[a_{i_0}^{-1}(X_n - a_0)]\right) = V(X_n).$$
        Hence, the base case holds. Next, suppose that the statement holds for some $r = k$ and consider the case where $V = V(f_1, \dots, f_k, f_{k+1})$. By the induction hypothesis, there is a polynomial $T$ such that
        $$V^T = V(f_1^T, \dots, f_k^T)\cap V(f_{k+1}^T) = V(X_{m+1}, ..., X_n, f_{k+1}^T).$$
        Now, notice that if we write $f = a_0 + \sum_{i=1}^{n}a_i X_i$, then $(X_{m+1}, ..., X_n, f_{k+1}^T) = (X_{m+1}, ..., X_n, g)$ where $g = a_0 + \sum_{i=1}^{m}a_i X_i$. Hence: $V^T = V(g, X_{m+1}, ..., X_n)$. If $g = 0$, we are done, if $g \neq 0$, then there is a $i = i_0 \leq m$ such that $a_{i_0} \neq 0$. Thus, if we define the polynomial map $U = (U_1, ..., U_n)$ by $U_i = 0$ for all $1\leq i\neq i_0 \leq m$, $U_{i_0} = a_{i_0}^{-1}(X_m - a_0)$, and $U_i = X_i$ for all $m+1 \leq i \leq n$. Thus, by construction:
        $$V^{T\circ U} = V(g\circ U, X_{m+1}\circ U, ..., X_n \circ U) = V(X_m, X_{m+1}, ..., X_n).$$
        Thus, the case $r = k+1$ also holds. Therefore, by induction, the proposition is proved for all $r \geq 1$.
        \item Suppose there exist affine change of coordinates $U_1$ and $U_2$ such that
        $$V^{U_1} = V(X_{m+1}, \dots, X_n) \qquad \text{ and } \qquad V^{U_2} = V(X_{s+1}, \dots, X_n)$$
        with $m < s$, then equivalently, there is a change of coordinates $T$ such that $V(X_{m+1}, ..., X_n)^T = V(X_{s+1}, ..., X_n)$. If we write $T = (T_1, ..., T_n)$, we get that $V(T_{m+1}, ..., T_n) = V(X_{s+1}, ..., X_n)$. By the Nullstellensatz, and the fact that $(X_{s+1}, ..., X_n)$ is a prime ideal:
        $$\{T_{m+1}, ..., T_n\} \subseteq \rad(T_{m+1}, ..., T_n) = I(V(T_{m+1}, ..., T_n)) = (X_{s+1}, ..., X_n).$$
        Thus, for all $m+1 \leq i \leq n$, $T_i \in (X_{s+1}, ..., X_n)$. Equivalently, 
        $$T_i = \alpha_{s+1, i}X_{s+1} + \dots + \alpha_{n,i}X_n$$
        for some scalars $\alpha_{j,i} \in k$, and for all $m+1 \leq i \leq n$ (technically, the coefficients should be polynomials but we can deduce that there must be such scalars by the fact that $T_i$ has degree 1). Now, recall that $T$ is an affine change of coordinates so its linear part must be invertible. Equivalently, the matrix form of its linear part must be invertible, and hence, its rows must be linearly independent vectors. But the above equation shows that the last $n -m$ rows of that matrix are of the form $(0,..., 0, \alpha_{s+1,i}, ..., \alpha_{n,i})$. In other words, the vectors
        $$\begin{pmatrix}
            0 \\ \vdots \\ 0 \\ \alpha_{s+1,m+1} \\ \vdots \\ \alpha_{n,m+1}
        \end{pmatrix}, \qquad \begin{pmatrix}
            0 \\ \vdots \\ 0 \\ \alpha_{s+1,m+2} \\ \vdots \\ \alpha_{n,m+2}
        \end{pmatrix}, \qquad \dots, \qquad \begin{pmatrix}
            0 \\ \vdots \\ 0 \\ \alpha_{s+1,n} \\ \vdots \\ \alpha_{n,n}
        \end{pmatrix}$$
        are linearly independent. This is equivalent to the fact that the vectors
        $$\begin{pmatrix}
            \alpha_{s+1,m+1} \\ \vdots \\ \alpha_{n,m+1}
        \end{pmatrix}, \qquad \begin{pmatrix}
            \alpha_{s+1,m+2} \\ \vdots \\ \alpha_{n,m+2}
        \end{pmatrix}, \qquad \dots, \qquad \begin{pmatrix}
            \alpha_{s+1,n} \\ \vdots \\ \alpha_{n,n}
        \end{pmatrix}$$
        are linearly. But this is impossible since we have $n - m$ vectors in $k^{n - s}$, and $n - m > n-s$. Therefore, we found a contradiction that proves that $m$ is independent of the choice of $T$.
        
        We now need to prove that $V$ is isomorphic to $\A^m(k)$ as a variety. Let $T$ be an affine coordinates change such that $V^T = V_m := V(X_{m+1}, ..., X_n)$, then $T(V_m) = V$. It follows that $T$, restricted to $V_m$, is a bijection with image equal to $V$. Since the inverse of $T$ is also an affine change of coordinates, then it follows that $V$ and $V_m$ are isomorphic as variety. Thus, we now need to show that $V_m$ is isomorphic to $\A^m(k)$ as variety. Consider the polynomial map $U: \A^m(k) \to V_m$ defined by
        $$U(X_1, ..., X_m) = (X_1, ..., X_m, 0, ..., 0).$$
        Its inverse is the polynomial map $\pi : V_m \to \A^m(k)$ defined by
        $$\pi(X_1, ..., X_m, X_{m+1}, ..., X_n) = (X_1, ..., X_m).$$
        By construction, they are inverses of each other so $U$ is an isomorphism of variety. Thus, $V$ is isomorphic to $V_m$, which is in turn isomorphic to $\A^m(k)$ as a variety. \\
    \end{enumerate}
\end{solution}

\begin{exercise}
    Let $P = (a_1, ..., a_n)$, $Q = (b_1, ..., b_n)$ be distinct points in $\A^n$. The \textit{line} through $P$ and $Q$ is defined to be $\{a_1 + t(b_1 - a_1), ..., a_n + t(b_n - a_n) | t\in k\}$. (a) Show that if $L$ is the line through $P$ and $Q$, and $T$ is an affine change of coordinates, then $T(L)$ is the line through $T(P)$ and $T(Q)$. (b) Show that a line is a linear subvariety of dimension 1, and that a linear subvariety of dimension 1 is the line through any of its two points. (c) Show that, in $\A^2$, a line is the same thing as a hyperplane. (d) Let $P,P' \in \A^2$, $L_1,L_2$ two distinct lines through $P$, $L_1', L_2'$ are distinct lines through $P'$. Show that there is an affine change of coordinates $T$ of $\A^2$ such that $T(P) = P'$ and $T(L_i) = L_i'$, $i=1,2$. \\
\end{exercise}

\begin{solution}
    \begin{enumerate}
        \item First, rewrite the line through $P$ and $Q$ as $L = \{P + t(Q-P) | t \in k\}$. If we let $T = T' + \alpha$ be an affine change of coordinates where $T'$ is the linear part, and $\alpha \in k^n$ is a constant, then 
        \begin{align*}
            T(L) &= \{T(P + t(Q-P)) | t \in k\} \\
            &= \{T'(P + t(Q-P)) + \alpha | t \in k\} \\
            &= \{T'(P) + t(T'(Q) - T'(P)) + \alpha | t \in k\} \\
            &= \{[T'(P) + \alpha] + t([T'(Q) + \alpha] - [T'(P) + \alpha]) | t \in k\} \\
            &= \{T(P) + t(T(Q) - T(P)) | t \in k\}
        \end{align*}
        which is precisely the definition of the line through $T(P)$ and $T(Q)$.
        \item First, consider the affine change of coordinates $T_1$ that maps $x$ to $x - P$ (a translation), then the image of $L$ is the line $L_1$ through 0 with director vector $v_1 = Q-P$. Next, since $v_1$ is a non-zero vector, then we can find vectors $v_2, ..., v_n$ such that $v_1, ..., v_n$ forms a basis of $k^n$. From that, define the linear transformation $T_2$ that maps $v_i$ to $e_i$ where $e_1, ..., e_n$ is the standard basis. By construction, $T_2$ is an invertible linear transformation, so it is an affine change of coordinates. The image of $L_1$ under $T_2$ is the line $L_2 = V(X_2, ..., X_n)$, which is a linear subvariety. Thus, by construction, we have that $L = T_1^{-1}(T_2^{-1}(V)) = V^{T_2 \circ T_1}$ is also a linear subvariety (by Problem 2.14 part (a)).
        
        Conversely, let $L$ be a linear subvariety of dimension 1, then there is an affine change of coordinates $T$, such that $L^T = V(X_2, ..., X_n)$, or equivalently, such that $L = T(V(X_2, ..., X_n))$. Let $P$ and $Q$ be two points in $L$, then there are two points $(p, 0, ..., 0), (q, 0, ..., 0) \in V(X_2, ..., X_n)$ that are the respective preimages of these two points. Since $P$ and $Q$ are distinct, and $T$ is a bijection, then $p$ and $q$ are also distinct. Hence, we can easily see that $V(X_2, ..., X_n)$ can be described as the line $L_0$ through the vectors $(p, 0, ..., 0)$ and  $(q, 0, ..., 0)$. Finally, using part (a), we get that $L = T(L_0)$ must be the line through the points $P$ and $Q$. Since $P \neq Q$ were arbitrary points on $L$, then we can say that a linear subvariety fo dimension 1 is the line through any two of its points.
        \item Let $L$ be a line in $\A^2$, then by the previous part, we know that there exists an affine change of coordinates $T$ such that $L^T = V(X_2)$, or equivalently, such that $L = T(V(X_2))$. Since $T$ is a bijection such that its inverse is also an affine change of coordinates, we can rewrite the previous equation as
        $$L = T(V(X_2)) = (T^{-1})^{-1}(V(X^2)) = V(X_2)^{T^{-1}}.$$
        Hence, if we write $T^{-1} = (T_1, T_2)$, we get that $L = V(X_2 \circ T^{-1}) = V(T_2)$. Therefore, $L$ is a hyperplane ($T_2$ has degree 1). 

        Conversely, let $V(f)$ be a hyperplane, then $f$ is a polynomial of degree 1. By the proof of the base case of the induction in the solution of part (b) in Problem 2.14, we have that there exists an affine change of coordinates $T$ such that $V(f)^T = V(X_2)$. Since the latter is the line through $(0,0)$ and $(1,0)$, and $T(V(X_2)) = V(f)$, then $V(f)$ must be a line. Therefore, in $\A^2$, hyperplanes and lines form the same class of objects.
        \item First, define the affine change of coordinates $T_1$ that acts as a translation by mapping the point $P$ to the origin. Thus, the image of the lines $L_1$ and $L_2$ are now intersecting at the origin. Let $v_1$ and $v_2$ be director vectors of the lines $T(L_1)$ and $T(L_2)$, since the lines are distinct, the vectors form a basis of $\A^2$. Similarly, repeat the same process for the lines $L_1'$ and $L_2'$: let $T_2$ be an affine change of coordinates that maps $P'$ to the origin (a translation), and let $v_3$, $v_4$ be the director vectors of the lines $T_2(L_1')$ and $T_2(L_2')$ respectively. Again, $v_3$, $V_4$ forms another basis of $\A^2$. Finally, consider the linear transformation that maps $v_1$ to $v_3$ and $v_2$ to $v_4$, then it is invertible, and so it is an affine change of coordinates. By construction, $T(T_1(L_1)) = T_2(L_1')$ and $T(T_1(L_2)) = T_2(L_2')$. Therefore, the affine change of coordinates $U = T_2^{-1}\circ T \circ T_1$ has the property that $U(P) = P'$ and $U(L_i) = L_i'$ for $i = 1,2$. \\
    \end{enumerate}
\end{solution}

\begin{exercise}
    Let $k = \C$. Give $\A^n(\C) = \C^n$ the usual topology (obtained by identifying $\C$ with $\R^2$, and hence $\C^n$ with $\R^{2n}$). Recall that a topological space $X$ is \textit{path-connected} if for any $P,Q \in X$, there is a continuous mapping $\gamma : [0,1] \to X$ such that $\gamma(0) = P$ and $\gamma(1) = Q$. (a) Show that $\C \setminus S$ is path-connected for any finite set $S$. (b) Let $V$ be an algebraic set in $\A^n(\C)$. Show that $\A^n(\C) \setminus V$ is path-connected. (\textit{Hint:} If $P,Q \in \A^n(\C) \setminus V$, let $L$ be the line through $P$ and $Q$. Then $L \cap V$ is finite, and $L$ is isomorphic to $\A^1(\C)$.) \\
\end{exercise}

\begin{solution}
    \begin{enumerate}
        \item Let $S = \{s_1, ..., s_n\}$ and consider the line segment $L$ between $P$ and $Q$. If $L\cap S = \varnothing$, then it suffices to take $\gamma(t) = p + t(q-p)$. Otherwise, take any point $l \in L$ and consider the straightline $L_0$ through $l$ that is perpendicular to $L$. For all points $z \in L_0$, we can define the line segments $L_1^{(z)}$ between $P$ and $z$, $L_2^{(z)}$ between $z$ and $Q$. For two distinct $z_1, z_2 \in L_0$, we have that
        $$(L_1^{(z_1)} \cup L_2^{(z_1)})\cap (L_1^{(z_2)} \cup L_2^{(z_2)}) = \{P,Q\}.$$
        Thus, if we define $L_z = L_1^{(z)} \cup L_2^{(z)}$, then there must be a $z_0 \in L_0$ such that $L_z \cap S = \varnothing$ since there are infinitely many points on $L_0$ but only finitely many in $S$. From that, define 
        $$\gamma(t) = \begin{cases}
            p+2t(z_0 - p) & 0 \leq t < \frac{1}{2}, \\
            z_0+2t(q - z_0) & \frac{1}{2} \leq t \leq 1. \\
        \end{cases}.$$
        After easily verifying that $\gamma$ is continuous, we are done.
        \item Let $P,Q \in \A^n(\C) \setminus V$ and $L = \{P + t(Q-P) | t \in \C\}$ be the line through $P$ and $Q$. If we write $V = V(f_1, ..., f_r)$, then
        $$L\cap V = (L\cap V(f_1)) \cap V(f_2, ..., f_r) \subseteq L\cap V(f_1).$$
        If we look at $L \cap V(f_1)$ more closely, we get that it is the set of solutions of the polynomial $f_1$ such that the input is in the line $L$. If we define $g(t) = P + t(Q-P)$, we get that the elements are precisely the image under $g$ of the solutions of the function $f_1(g(t))$. But notice now that $f_1 \circ g$ is only one variable so it has finitely many solutions. Thus, $L \cap V(f_1)$ is finite, and hence, so is $L \cap V$.
        
        Next, since $L$ is a linear subvariety of dimension 1, then it is isomorphic to $\A^1(\C)$ as a variety (Problem 2.14 part (c)). Hence, there is a bijective polynomial map $T : \C \to L$. Let $S = T^{-1}(L \cap V)$, then $S$ is finite. By part (a), there is a path $\gamma : [0,1] \to \C\setminus S$ from $P' = T^{-1}(P)$ to $Q' = T^{-1}(Q)$. It follows that $T \circ \gamma : [0,1] \to L \setminus V \subseteq \A^n(\C)\setminus V$ is path from $P$ to $Q$. Therefore, $\A^n(\C)$ is path-connected. \\
    \end{enumerate}
\end{solution}